\chapter{Comparison with Nektar++}\label{ch:nektar}

The aim is to run all the Aconcagua cases with OpenFOAM and with Nektar++ and to
compare the two solvers. Nektar++ adds $p$-refinement to the $h$-refinement of the
meshes, so it has more cases.

``The same case'' is the same mesh, the same free stream and the same initial field.
The schemes, the time integration, the compressible formulation and the turbulence
treatment can differ, but they must be as near as possible. This chapter gives what
each OpenFOAM run records for the Nektar++ run, and what to measure in the same way in
the two solvers.

\section{What must be the same}\label{sec:nektar-same}

\begin{longtable}{@{}P{0.18\linewidth} P{0.34\linewidth} P{0.40\linewidth}@{}}
\caption{Items that the Nektar++ run must reproduce.}\label{tab:nektar-same}\\
\toprule
Item & Where it is & Note \\
\midrule
\endfirsthead
\toprule
Item & Where it is & Note \\
\midrule
\endhead
Mesh & \ent{meshHash} in \file{run.info}; \file{<mesh-dir>/meshHash} & sha256 (16 characters) of \file{points}, \file{faces}, \file{owner}, \file{neighbour}, \file{boundary}. With \ent{--refine}, the hash of the refined mesh. The Nektar++ mesh comes from the same \file{polyMesh}. \\
Half mesh & \file{meshInfo}: \ent{sector half} & Symmetry plane $z = 0$. \ent{Aref} is half the frontal area. \\
Free stream & block \ent{flow} in \file{run.info} & \ent{Minf}, \ent{alpha}, \ent{pInf}, \ent{Tinf}, \ent{Uinf}, \ent{rhoInf}, \ent{muInf}, the components of $\mathbf{U}$. \\
Reynolds number & \ent{RePerMetre}, \ent{ReD}, \ent{ReL} & From \ent{rhoInf}, \ent{Uinf} and \ent{muInf}, on the diameter and on the body length. \\
Gas & \ent{Rgas}, \ent{gammaGas}, \ent{Cp}, \ent{RgasThermo} & $R$ is 287.05 in \file{flowDerived} and 287.10 in the thermophysical model (\cref{sec:thermo}). \\
Viscosity & \ent{sutherland (As Ts)} & $\mu = A_s\sqrt{T}/(1 + T_s/T)$. \\
Conductivity & \ent{PrEucken} & $Pr = 0.690$ from the Eucken correction, not 0.72. Nektar++ must use this value. \\
Turbulence & \ent{turbulence}, \ent{Prt}, \ent{Ti}, \ent{nuRatio}, \ent{kInlet}, \ent{omegaInlet} & $k$--$\omega$ SST with \ent{nutUSpaldingWallFunction}. \\
Wall & \ent{wallT}, \ent{wallNut} & Adiabatic (\val{zeroGradient} on $T$), no slip. \\
Initial field & \ent{initial} & \val{uniform free stream}, or the run that was mapped. \\
References & \file{FOcommon} & \ent{CofR (0 0 0)} (nose tip); drag and lift in wind axes; \ent{lRef}; \ent{Aref}. \\
\bottomrule
\end{longtable}

\section{What to compare}

\begin{tabularx}{\linewidth}{@{}P{0.32\linewidth} L@{}}
\toprule
Quantity & OpenFOAM file \\
\midrule
History of $C_D$, $C_L$, $C_m$ & \file{forceCoeffs1/*/coefficient.dat}, every 50 steps \\
Pressure and viscous force & \file{forces1/*/force.dat} \\
$C_p$, $\tau_w$, \yplus{} along the body & \file{samples/*/skinZ0_*.csv}, \file{skinY0_*.csv} \\
Wake profiles & \file{samples/*/wake05}, \file{wake1}, \file{wake2}, \file{wake5} \\
Axis and free-stream check & \file{samples/*/axis_*.csv}, \file{freeStream_*.csv} \\
Skin and symmetry plane & \file{surfaces/*/} (VTK) \\
Cost & \ent{cyclesPerCellStep}, \ent{cpuSeconds}, \ent{energyPkgJ}, \ent{GHzEffective} in \file{run.info} \\
\bottomrule
\end{tabularx}

\section{How to measure in the same way}

\begin{itemize}
\item \textbf{Degrees of freedom.} OpenFOAM: $N_\text{cells}$ for each variable.
  Nektar++: $N_\text{elements}\,(P+1)^3$ for each variable on hexahedra (a different
  count on prisms and tetrahedra). Compare the cost for each degree of freedom and
  each step, in cycles.
\item \textbf{Convergence.} Each solver stops with its own criterion and its own
  step. For the cost to the result, apply one criterion to the two histories of $C_D$,
  against the accumulated cost. Do not use the residuals: each code normalizes them
  differently.
\item \textbf{Refinement.} Each pair $(h, P)$ of Nektar++ is a case. The cost for each
  degree of freedom puts them on the same curve as $h1$ to $h3$ of OpenFOAM.
\item \textbf{Computer.} Use the procedure of \cref{sec:measurement}, the same
  computer, the same number of processes and the same binding. Do not divide by the
  clock frequency: measure the dependence on the clock.
\end{itemize}
